% !TEX program = xelatex
% 학회 참관 리서치 리포트 템플릿 — <conf>/report/report.tex 로 복사해 사용.
% 검증된 원본: icml_2026/report/report.tex (ICML 2026)
% 빌드: make  (같은 디렉토리의 references/Makefile 복사)
\documentclass[10pt]{article}

% ---------- 언어/폰트 ----------
\usepackage{kotex}
\usepackage{fontspec}
\setmainhangulfont{Apple SD Gothic Neo}
\setsanshangulfont{Apple SD Gothic Neo}

% ---------- 레이아웃 ----------
\usepackage{amsmath}
\usepackage[a4paper,margin=22mm,headsep=7mm]{geometry}
\usepackage{xcolor}
\usepackage{colortbl}
\usepackage{graphicx}
\usepackage{booktabs}
\usepackage{array}
\usepackage{multirow}
\usepackage{enumitem}
\usepackage{needspace}
\usepackage{caption}
\usepackage{titlesec}
\usepackage{fancyhdr}
\usepackage[hidelinks]{hyperref}

% EXIF 회전을 픽셀에 반영한 사본(assets/)을 우선 사용 — 원본은 fallback
% (assets는 scripts/normalize_photos.py 로 생성)
\graphicspath{{assets/research/}{assets/}{assets/booth/}{../pictures/research/}{../pictures/}{../pictures/booth/}}

% ---------- 색/스타일 ----------
\definecolor{icmlblue}{HTML}{1F3A93}
\definecolor{accent}{HTML}{2E86C1}
\definecolor{softgray}{HTML}{6B7280}
\definecolor{rule}{HTML}{D5DBDB}
\definecolor{labelbg}{HTML}{EEF1F8}   % 스펙 표 라벨열 배경
\definecolor{resultbg}{HTML}{EAF3FA}  % 핵심 결과 박스 배경

\hypersetup{colorlinks=true, linkcolor=icmlblue, urlcolor=accent, citecolor=accent}
\captionsetup{font=small, labelfont={bf,color=icmlblue!85!black}}

\titleformat{\section}{\large\bfseries\color{icmlblue}}{\thesection}{0.6em}{}[{\color{rule}\titlerule[0.8pt]}]
\titlespacing*{\section}{0pt}{14pt}{7pt}

\setlist[itemize]{leftmargin=1.4em, topsep=2pt, itemsep=1.5pt}

\pagestyle{fancy}
\fancyhf{}
\fancyhead[L]{\small\color{softgray} <학회명> 참관 리서치 리포트}
\fancyhead[R]{\small\color{softgray} <개최지>}
\fancyfoot[C]{\thepage}
\renewcommand{\headrulewidth}{0.4pt}
\renewcommand{\footrulewidth}{0pt}

% 긴 URL이 여백을 넘지 않도록 줄바꿈 허용
\Urlmuskip=0mu plus 1mu\relax
\def\UrlBreaks{\do\/\do\-\do\.\do\_\do\a\do\b\do\c\do\d\do\e\do\f\do\g\do\h\do\i\do\j\do\k\do\l\do\m\do\n\do\o\do\p\do\q\do\r\do\s\do\t\do\u\do\v\do\w\do\x\do\y\do\z}
\emergencystretch=2em

% ---------- 헬퍼 매크로 ----------
% 섹션 제목이 페이지 하단에 홀로 남지 않도록 최소 공간 확보
\newcommand{\sect}[1]{\needspace{5\baselineskip}\section{#1}}
% 오프닝 갤러리 사진 (2열, 높이 통일 — 세로/가로 혼합에도 행이 균일)
\newcommand{\gphoto}[2]{%
  \begin{minipage}[t]{0.485\linewidth}\centering
  \includegraphics[height=48mm,keepaspectratio]{#1}\par\vspace{2pt}
  {\footnotesize\color{softgray}#2}\end{minipage}}
% 논문 단위 통짜 박스(minipage) — 페이지 경계에서 잘리지 않고, 번호 매긴 소제목을 목차에 등록
\newenvironment{paperbox}[1]{%
  \par\addvspace{6pt}%
  \refstepcounter{subsection}%
  \phantomsection
  \addcontentsline{toc}{subsection}{\protect\numberline{\thesubsection}#1}%
  \noindent\begin{minipage}{\linewidth}%
  \setlength{\parindent}{0pt}%
  {\normalsize\bfseries\color{icmlblue!85!black}\thesubsection\quad #1\par}\vspace{3pt}%
}{%
  \end{minipage}\par\addvspace{10pt}%
}
\newcommand{\meta}[1]{{\footnotesize\color{softgray}\sffamily\raggedright #1\par}\vspace{3pt}}
% 논문 스펙 카드: 연구 배경 / 기존 방법의 한계 / 정의한 문제 / 해결 방법
% 주의: \columncolor는 \multirow와 함께 쓰면 라벨이 가려진다 — multirow 표에는 쓰지 말 것.
\newcommand{\pcard}[4]{%
{\footnotesize\renewcommand{\arraystretch}{1.3}%
\begin{tabular}{@{}>{\columncolor{labelbg}\bfseries\color{icmlblue!88!black}}p{2.35cm} p{13.5cm}@{}}
\toprule
연구 배경 & #1\\ \midrule
기존 방법의 한계 & #2\\ \midrule
정의한 문제 & #3\\ \midrule
해결 방법 & #4\\
\bottomrule
\end{tabular}}\par\vspace{4pt}}
% 핵심 결과: 옅은 배경 박스로 시각 분리
\newcommand{\result}[1]{%
{\setlength{\fboxsep}{4pt}%
\noindent\colorbox{resultbg}{\parbox{\dimexpr\linewidth-10pt\relax}{\footnotesize\textbf{\color{accent}핵심 결과.}\ #1}}\par}\vspace{4pt}}
% 그림 (paperbox 내부) — 중앙 정렬 이미지 + 캡션. float 금지(잘림 방지).
\newcommand{\shot}[3]{\par\vspace{2pt}{\centering\includegraphics[width=#2\linewidth]{#1}\par\vspace{1pt}\captionof{figure}{#3}\par}}

% ================================================================
\begin{document}

% ---------- 표지 ----------
\begin{titlepage}
\centering
\vspace*{3.0cm}
{\color{softgray}\large <학회 정식 명칭>}\\[4pt]
{\color{icmlblue}\Huge\bfseries <학회명> 참관 리서치 리포트}\\[10pt]
{\large 포스터 · 구두 발표 리뷰}\\[2.2cm]
{\large 주제 : <주제 목록>}\\[6pt]
{\large 정리 대상 : 인상 깊었던 논문 · 발표 N편}\\[2.2cm]
{\color{softgray} 개최지 : <도시> (<연월>)}\\[3pt]
{\color{softgray} 작성일 : <연월>}\\
\vfill
\end{titlepage}

% ---------- 오프닝: 현장 스케치 (2열 갤러리) ----------
\newpage
\thispagestyle{fancy}
{\centering
{\Large\bfseries\color{icmlblue} <학회명> · <도시> — 현장 스케치}\par\vspace{3pt}
{\small\color{softgray} 등록 · 학회장 · Sponsor Expo}\par}
\vspace{12pt}

\noindent\gphoto{<사진1>}{<캡션1>}\hfill\gphoto{<사진2>}{<캡션2>}\par\vspace{10pt}
% ... 행 반복 (페이지당 4행 권장)

% ---------- 목차 ----------
\newpage
\renewcommand{\contentsname}{\color{icmlblue}Contents \normalfont\normalsize (목차)}
\tableofcontents
\thispagestyle{fancy}

% ================================================================
% 탑다운 구성: §1 연구 동향(표) → §2 개요·전체 목록(표) → §3~ 주제별 상세
\newpage
\section{연구 동향 종합 (Overall Research Trends)}

% 동향 표: >{\columncolor{labelbg}...}p{3.0cm} p{9.4cm} p{3.2cm}, 셀 핵심 키워드 \textbf
% + 미리뷰 Spotlight 표(약칭 \textbf — 한줄 요약, abstract 기반) — minipage 통짜

\sect{개요 및 전체 목록 (Overview)}

% 14편 목록 표 (분야/제목/소속/형태) — multirow 사용 시 columncolor 금지

\sect{<주제 1>}

% 섹션마다 도입 1~2문장 필수 (제목 고아 방지 + 맥락)

\begin{paperbox}{<약칭> — <한줄 제목>}
\meta{<저자> · <소속> · <형태: 포스터/Spotlight (oral talk)> · arXiv:<id>}
\pcard
{<연구 배경 — 도메인 맥락 1~2문장, 핵심 키워드 \textbf>}
{<기존 방법의 한계 — 실패 지점 \textbf>}
{<정의한 문제 — 문제 정식화, 핵심 개념 \textbf>}
{<해결 방법 — 방법명 \textbf, 파이프라인은 $\rightarrow$ 구조>}
\result{<검증된 수치 중심 1~3문장, 대표 수치 \textbf>}
\shot{<사진파일>}{0.42}{<캡션>}
\end{paperbox}

\end{document}
